% ===== PRÉAMBULE CANONIQUE — à coller VERBATIM en tête de CHAQUE .tex =====
\documentclass[11pt,a4paper]{article}
\usepackage[utf8]{inputenc}\usepackage[T1]{fontenc}\usepackage{lmodern}
\usepackage[french]{babel}
\usepackage{amsmath,amssymb,mathtools}\usepackage{icomma}
\usepackage[version=4]{mhchem}          % \ce{...} : équations & formules
\usepackage{siunitx}\sisetup{locale=FR} % \qty{}{}, \si{}, \ang{}
\usepackage{chemfig}                    % \chemfig{...} : structures développées
\usepackage{xcolor}\usepackage{colortbl}
\usepackage{tikz}\usepackage{pgfplots}\pgfplotsset{compat=1.18}
\usepackage[most]{tcolorbox}\usepackage{enumitem}\usepackage{array,booktabs}
\usepackage[a4paper,margin=2.2cm]{geometry}
\usetikzlibrary{arrows.meta,calc,angles,quotes,positioning,patterns,backgrounds,shapes.geometric}
\definecolor{primary}{RGB}{222,49,99}\definecolor{primarydark}{RGB}{150,22,55}
\definecolor{defcol}{RGB}{0,114,178}\definecolor{methcol}{RGB}{52,92,148}
\definecolor{excol}{RGB}{0,135,90}\definecolor{attncol}{RGB}{200,40,55}
\tcbset{boxbase/.style={enhanced,breakable,boxrule=0.9pt,arc=2.5pt,
  left=3mm,right=3mm,top=2.3mm,bottom=2.3mm,fonttitle=\bfseries,coltitle=white}}
% --- boîtes TUTORIEL (noms EXACTS, ne pas renommer) ---
\newtcolorbox{cle}[1][]{boxbase,colback=defcol!6,colframe=defcol,
  title={\textsc{À retenir}\ifx\relax#1\relax\relax\else\textmd{ : \itshape #1}\fi}}
\newtcolorbox{methode}[1][]{boxbase,colback=methcol!7,colframe=methcol,sharp corners,
  title={\textsc{Méthode}\ifx\relax#1\relax\relax\else\textmd{ --- \itshape #1}\fi}}
\newtcolorbox{exemplebox}[1][]{boxbase,colback=excol!5,colframe=excol,
  title={\textsc{Exemple}\ifx\relax#1\relax\relax\else\textmd{ : \itshape #1}\fi}}
\newtcolorbox{piege}[1][]{boxbase,colback=attncol!6,colframe=attncol,
  title={\textsc{Piège}\ifx\relax#1\relax\relax\else\textmd{ : \itshape #1}\fi}}
\newtcolorbox{remarque}[1][]{boxbase,colback=black!4,colframe=black!45,
  title={\textsc{Remarque}\ifx\relax#1\relax\relax\else\textmd{ : \itshape #1}\fi}}
% --- boîtes EXERCICE / CORRIGÉ (noms EXACTS, auto-comptées) ---
\newtcolorbox[auto counter]{exo}[1][]{boxbase,colback=primary!4,colframe=primary,
  title={\textsc{Exercice}~\thetcbcounter\ifx\relax#1\relax\relax\else\textmd{ --- \itshape #1}\fi}}
\newtcolorbox[auto counter]{corrige}[1][]{boxbase,colback=excol!4,colframe=excol,
  title={\textsc{Corrigé}~\thetcbcounter\ifx\relax#1\relax\relax\else\textmd{ --- \itshape #1}\fi}}
\newcommand{\R}{\mathbb{R}}\newcommand{\N}{\mathbb{N}}\newcommand{\Z}{\mathbb{Z}}\newcommand{\ds}{\displaystyle}
% (un \usepackage{fancyhdr}+en-tête est possible ; il est ignoré côté web)
% =========================================================================
\begin{document}
\sisetup{per-mode=symbol}

\begin{corrige}[choisir et doser un couple tampon]
\begin{enumerate}[label=\alph*)]
  \item Un tampon est efficace pour $\mathrm{pH} \in [\mathrm{p}K_a - 1,\ \mathrm{p}K_a + 1]$. Pour
        $\mathrm{pH} = 5{,}0$, le couple dont le $\mathrm{p}K_a$ est le plus proche est
        \textbf{\ce{CH3COOH/CH3COO-}} ($\mathrm{p}K_a = 4{,}8$).
  \item $\dfrac{[\ce{A-}]}{[\ce{HA}]} = 10^{\,\mathrm{pH}-\mathrm{p}K_a} = 10^{\,5{,}0-4{,}8} = 10^{0{,}2}
        \approx \mathbf{1{,}6}$ (car $\log 1{,}6 = 4\log 2 - 1 = 0{,}20$).
\end{enumerate}
\end{corrige}

\begin{corrige}[le tampon à la demi-équivalence]
\begin{enumerate}[label=\alph*)]
  \item Le pH vaut $\mathrm{p}K_a$ quand $[\ce{HA}] = [\ce{A-}]$, c'est-à-dire quand on a neutralisé la
        \textbf{moitié} de l'acide. Quantité initiale : $n_0 = 0{,}100 \times 0{,}20 = \qty{0,020}{mol}$ ; il faut
        $n(\ce{OH-}) = 0{,}010$ mol, soit
        \[ V(\ce{NaOH}) = \frac{0{,}010}{0{,}20} = \qty{0,050}{L} = \mathbf{50\ mL}. \]
  \item C'est la \textbf{demi-équivalence} du titrage.
\end{enumerate}
\end{corrige}

\begin{corrige}[titrage par un diacide]
\begin{enumerate}[label=\alph*)]
  \item Chaque \ce{H2SO4} libère \textbf{deux} \ce{H+}, donc à l'équivalence $\;2\,C_a V_a = C_b V_b$.
  \item $V_a = \dfrac{C_b V_b}{2\,C_a} = \dfrac{0{,}10 \times 30}{2 \times 0{,}050} = \dfrac{3{,}0}{0{,}10}
        = \mathbf{30\ mL}$.
\end{enumerate}
\end{corrige}

\begin{corrige}[un tampon encaisse un ajout d'acide fort]
\begin{enumerate}[label=\alph*)]
  \item L'acide fort est consommé par la base du tampon : $\ce{HCOO- + H3O+ -> HCOOH + H2O}$. Quantités
        initiales : $n(\ce{HCOOH}) = n(\ce{HCOO-}) = 0{,}100\times0{,}50 = \qty{0,050}{mol}$ ;
        $n(\ce{H3O+}) = 0{,}010\times1{,}0 = \qty{0,010}{mol}$. Après réaction :
        \[ n(\ce{HCOO-}) = 0{,}050 - 0{,}010 = \qty{0,040}{mol},\qquad n(\ce{HCOOH}) = 0{,}050 + 0{,}010 = \qty{0,060}{mol}. \]
  \item $\mathrm{pH} = 3{,}8 + \log\dfrac{0{,}040}{0{,}060} = 3{,}8 + \log\tfrac{2}{3}
        = 3{,}8 + (0{,}30 - 0{,}48) = 3{,}8 - 0{,}18 = \mathbf{3{,}62}$. Le pH n'a presque pas bougé : le tampon a
        bien encaissé l'ajout.
\end{enumerate}
\end{corrige}

\begin{corrige}[quel sel donne une solution basique ?]
\begin{enumerate}[label=\alph*)]
  \item \ce{Ca(NO3)2} : \ce{Ca(OH)2} (base forte) + \ce{HNO3} (acide fort) ;\quad
        \ce{KCN} : \ce{KOH} (base forte) + \ce{HCN} (acide faible) ;\quad
        \ce{NaCl} : \ce{NaOH} (forte) + \ce{HCl} (fort) ;\quad
        \ce{NH4Br} : \ce{NH3} (base faible) + \ce{HBr} (acide fort).
  \item Le sel basique est \textbf{\ce{KCN}} (acide faible + base forte) : \ce{CN- + H2O <=> HCN + OH-}.
        (\ce{Ca(NO3)2} et \ce{NaCl} sont neutres ; \ce{NH4Br} est acide.)
\end{enumerate}
\end{corrige}

\end{document}
